\subsection{The global hypothesis in C.28: checking the prerequisite rather than assuming it}
\begin{lemma}[The no-small-quotients resolution used by INDLOCAL]\label{F:large-resolution}
For a connected CW complex $W$ and an infinite cardinal $\lambda$ at
least its number of cells, there is an acyclic map $t:BG\to W$ with
$|G|\leq\lambda$ such that every homomorphism from $G$ to a group of
cardinality less than $\lambda$ kills $\ker\pi_1(t)$. Here acyclic means
an isomorphism in homology with every local coefficient system on $W$.
\end{lemma}
\begin{proof}
We check the supplied \textsc{OWN\_RESULTS} Theorem~3.1, including its
previously unopened input. Leary's Theorem~A \cite[p.~2]{leary2013}
gives an aspherical locally CAT(0) complex mapping to any simplicial
complex by a homology equivalence for all local systems. A CW complex
has a simplicial model of cardinality at most the maximum of its cell
cardinality and $\aleph_0$. Leary's construction is finite on each finite
subcomplex and commutes with their unions. It therefore has that same
cardinal bound, and its fundamental group does too. This supplies the
exact Kan--Thurston input in the successor step; the original 1976 paper
need not be accepted from memory.

\textbf{The large simple acyclic overgroup input.} The embedding
construction checked in Theorem~\ref{A:acyclic} gives an acyclic
overgroup of any group without increasing an infinite cardinal bound.
For completeness the simple-envelope step used in the supplied argument
can be obtained as follows. For every $g\ne1$ and $h$ in a group $L$,
introduce a free letter $t=t_{g,h}$ and put
\[
 w=(gtg^{-1}t^{-1})(gt^2g^{-1}t^{-2}).
\]
Both $w$ and $wh$ have infinite order in $L*F(t_{g,h})$: their cyclically
reduced normal forms have positive alternating length (length eight,
or six after the possible boundary cancellation). Adjoin an HNN letter
$s$ with $s^{-1}ws=wh$. The cyclic subgroups being identified are
infinite cyclic, so the normal-form theorem embeds the old group.
Moreover $h=w^{-1}s^{-1}ws$ belongs to the normal closure of $g$.
Iterating over all pairs and then over countably many stages produces
an overgroup in which every nonidentity element normally generates the
whole group. Start with a free product with a free group on $\lambda$
letters to ensure cardinality exactly $\lambda$. Alternately take this
simple envelope and an acyclic envelope. In the union, every homology
class dies at a later acyclic stage, and every nonidentity element
normally generates every later stage. Thus the union is simple,
acyclic, of cardinality $\lambda$, and contains the specified group.
This checks the exact strength of CATHOT B.7 needed here.

\textbf{Start of the cellular induction.} Use a single vertex for $W$
and take $G_1$ to be the free group of its one-skeleton. Then
$BG_1\simeq W^{(1)}$, with zero kernel. Well-order higher cells so their
boundaries occur earlier; at limit stages use unions.

\textbf{The comparison for one attached cell.} Suppose $t_\alpha:BG_\alpha
\to W_\alpha$ has been constructed and the next cell has attaching map
$a:S^m\to W_\alpha$, $m\geq1$. For $m=1$, take $H'=\Z$ and lift the
attaching loop using the surjection on fundamental groups. For $m\geq2$,
form $E=S^m\times^h_{W_\alpha}BG_\alpha$. Acyclicity is preserved by
homotopy pullback, so $E\to S^m$ is acyclic and $E$ is connected.
Take the preceding bounded Kan--Thurston resolution $BH'\to E$.
The projection to $BG_\alpha$ is represented by a homomorphism
$u:H'\to G_\alpha$, and there is a specified homotopy
$t_\alpha Bu\simeq aq$ with $q:BH'\to S^m$ acyclic. The pullback can be
modeled with at most $\lambda$ simplices: finite words, horns, and
countably many simplicial replacement steps preserve an infinite
cardinal bound. Hence $|H'|\leq\lambda$.

\textbf{The amalgam and its map.} Embed $H'$ in two simple acyclic
groups $A',A''$ of size $\lambda$ and set
\[
 G_{\alpha+1}=(G_\alpha\times A')*_{H'}A'',
 \qquad h\longmapsto(u(h),i'(h)),\quad h\longmapsto i''(h).
\]
Both embeddings are injective, so both factors and $G_\alpha$ embed.
The Bass--Serre tree shows that the double mapping cylinder of their
classifying spaces is $BG_{\alpha+1}$: its universal cover is the tree
of contractible vertex and edge spaces. Map the first vertex space by
$t_\alpha\operatorname{pr}_1$, the second to the center of the new cell,
and the intervening cylinders by the chosen homotopy and the radial
cone on $q$. This defines $t_{\alpha+1}$ with the correct restriction.

\textbf{Acyclicity and the quotient property.} With a local system from
$W_{\alpha+1}$, the $A'$ factor acts trivially and is acyclic. The product
projection is consequently a local-coefficient homology equivalence;
the other two pieces are handled by acyclicity of $q$ and $A''$.
Mayer--Vietoris gives acyclicity of $t_{\alpha+1}$. A homomorphism into a
group smaller than $\lambda$ kills both simple factors. The amalgamation
relation then kills $u(H')$. By induction it factors through
$\pi_1W_\alpha$, and it kills the attaching loop when $m=1$; for $m\geq2$
there is no new fundamental-group relation. Van Kampen therefore makes
it factor through $\pi_1W_{\alpha+1}$. These are the comparison maps
needed in the supplied proof, not merely an equality of homology groups.

\textbf{Limits.} At a limit ordinal every finite chain and every group
word occurs earlier. Local-coefficient homology and the presentation of
the fundamental group commute with these unions. The factor maps are
unique because $G_\alpha\to\pi_1W_\alpha$ is onto, so they assemble.
The number of stages and every stage size are at most $\lambda$.
This completes all induction and cardinal steps of the asserted resolution.
\end{proof}

\begin{theorem}[The global ind-groupoid comparison used in C.28]\label{F:global-ind-groupoids}
Let $K:I\to\Cat_{\mathrm{hot}}$ be a small filtered diagram. If
$H(X)=\colim_i\pi_0\operatorname{Fun}(X,K_i)$ inverts every nerve
weak equivalence between small categories, then the canonical
ind-morphism $(K_i)\to(\Pi_1K_i)$ is an isomorphism. The global
quantifier on tests is essential.
\end{theorem}
\begin{proof}
\textbf{Descent and its specified map.} Invariance descends $H$ to a
presheaf $F$ on $\mathrm{Hot}$. For any presheaf $T$ there, Yoneda gives
\[
 \operatorname{Nat}(F,T)
 =\operatorname{Nat}(H,T\circ B)
 =\lim_i T(BK_i)
 =\operatorname{Nat}\bigl(\colim_i[-,BK_i],T\bigr).
\]
Thus $F(Y)=\colim_i[Y,BK_i]$, and the identification sends a functor to
its classifying map. The localization fact used here follows by replacing
a simplicial set by its category of simplices: its nerve realizes to the
original space, the last-vertex functor compares it to an original nerve,
and these natural equivalences represent maps and homotopies by roofs.
This also checks the naturality, which the subsequent cancellation needs.

\textbf{The two comparisons.} Choose an infinite cardinal $\mu$ bounding
all arrow sets of the stages, and choose $\lambda>\mu$ large enough for
a test space $Y$. Componentwise apply Lemma~\ref{F:large-resolution} to
obtain $t:T=\coprod BG_c\to Y$ with $p:G_c\twoheadrightarrow\pi_1Y_c$.
Every functor $\mathbf BG_c\to K_i$ lands in the core groupoid of $K_i$.
Its image automorphism group has at most $\mu$ elements, so the functor
factors uniquely through $\mathbf B\pi_1Y_c$. The same assertion for a
natural transformation follows by choosing a lift of each group element
under $p$ in the naturality equation. Hence
\[
 \operatorname{Fun}\bigl(\coprod\mathbf B\pi_1Y_c,K_i\bigr)
 \cong \operatorname{Fun}\bigl(\coprod\mathbf BG_c,K_i\bigr).
\]
This is the actual functor-category comparison in INDLOCAL B.1.

\textbf{The injectivity needed for cancellation.} An acyclic map
$t:T\to Y$ is onto on fundamental groups. The fold
$Y\cup_T^hY\to Y$ is an isomorphism on $\pi_1$ by van Kampen and on
homology with every local coefficient system by Mayer--Vietoris.
Using coefficients $\Z[\pi_1Y]$ gives a homology equivalence of simply
connected universal covers, hence a homotopy equivalence. Its two
sections are equal in $\mathrm{Hot}$. It follows that
$[Y,Z]\to[T,Z]$ is injective for every space $Z$. Filtered colimits
preserve this injectivity: an equality is witnessed at a common stage,
where the preceding injection applies.

\textbf{Cancellation and the unit.} Let $v:Y\to P_1Y=\coprod B\pi_1Y_c$.
The composites $vt$ and $Bp$ agree because their maps on fundamental
groupoids agree and the target is a $1$-type. The functor-category
comparison proves $F(vt)$ bijective. The previous paragraph proves
$F(t)$ injective. Since $F(vt)=F(t)F(v)$, surjectivity of the composite
also makes $F(t)$ surjective; therefore both $F(t)$ and $F(v)$ are
bijections. In particular $H$ inverts every categorical unit
$X\to\Pi_1X$. Yoneda now identifies the unit of the induced reflection
on ind-objects with $(K_i)\to(\Pi_1K_i)$: left Kan extension of each
representable is represented by its groupoid completion, and the unit
on their colimit is invertible exactly when $H$ inverts these units.
This verifies the filtered proof of INDLOCAL B.1 used in C.28; it does
not apply this theorem to a presheaf invariant only on finite tests.
\end{proof}

\begin{proposition}[C.28: the finite-test comparison under global invariance]
\label{F:C28}
Suppose a sequence $(K_n)$ has globally invariant ind-presheaf $H$ as
above. Put $Z=\hocolim_n BK_n$ and
$\Gamma=\colim_n\Pi_1K_n$. The canonical map $Z\to B\Gamma$ is a
homotopy equivalence. For every finite category $Y$ the specified map
\[
 c_Y:H(Y)\longrightarrow[BY,Z]
\]
is a bijection. The same conclusion holds when $\Pi_1Y$ is finitely
presented. This is the valid finite-category formula at VII
ms.~129--132, under the global invariance premise of that paragraph.
\end{proposition}
\begin{proof}
\textbf{The map of spaces.} The stagewise units give the displayed
canonical map. Finite edge words, relations, and basepoint paths occur
at a common stage, so it is bijective on components and an isomorphism
on every fundamental group. The preceding theorem says $H$ depends
only on $\Pi_1$. A finite face poset for $S^r$, $r\geq2$, has trivial
fundamental groupoid, so its value equals the value at a point.
The descended formula in the proof above and compactness of the finite
sphere identify this with $[S^r,Z]=[*,Z]$. In each component this says
that every based homotopy class lies in the orbit of zero; that orbit is
just zero. Thus all higher homotopy groups vanish, and the canonical map
between these $1$-types is a homotopy equivalence.

\textbf{The finite-presentation step.} The groupoid $\Pi_1Y$ has a finite
presentation from objects, arrows, identity equations, composition
equations, and formal inverses. A functor to $\Gamma$ and a natural
isomorphism between two such functors are finite data with finitely many
equations. They therefore occur at a common stage. Consequently
\[
 H(Y)\cong\colim_n\pi_0\operatorname{Fun}(Y,\Pi_1K_n)
 \cong\pi_0\operatorname{Fun}(Y,\Gamma)
 \cong[BY,B\Gamma]\cong[BY,Z].
\]
Every arrow in this chain is induced by the stagewise units or their
classifying maps. Hence this is $c_Y$, not an unspecified bijection.
This checks both (C.28.1) and (C.28.2) and explains why no compactness
assumption about the whole nerve of $Y$ is needed.

\textbf{Why arbitrary tests are different.} Even the discrete groupoid
sequence $\{1,\ldots,n\}$ has a globally invariant ind-presheaf.
On an infinite discrete test set its colimit consists of finite-image
functions, whereas maps to its union $\N$ include functions with infinite
image. This is a fully explicit failure of the unrestricted version of
$c_Y$, with the same comparison map. Thus the qualification in C.28 is
necessary.
\end{proof}

\begin{lemma}[The required diagram-localization sphere test]
\label{F:diagram-sphere}
Let $\mathcal M$ be bounded finite-dimensional rational chain complexes,
marked at quasi-isomorphisms, and let $P$ be a finite poset. For constant
diagrams $A=\Q[0]$, $B=\Q[2]$, there is a specified isomorphism
\[
 \Hom_{\mathcal M^P[\Sigma_P^{-1}]}(A_P,B_P)
       \cong H^2(BP;\Q).
\]
In particular the values for a finite face poset of $S^2$ and for a
point are $\Q$ and the zero group, respectively.
\end{lemma}
\begin{proof}
Use the diagram of chain complexes
$Q_P(x)=C_*(N(P\downarrow x);\Q)$. Its augmentation to $A_P$ is an
objectwise quasi-isomorphism because each slice has terminal object
$x$. In degree $r$ it is a finite sum of projective representable
diagrams, one for each strict chain $x_0<\cdots<x_r$, with summand
$\Q[P(x_r,-)]$. Thus $Q_P$ is a bounded complex of finite projective
diagrams. The category of diagrams of bounded finite-dimensional
complexes on finite $P$ is exactly the category of bounded complexes
of finite-dimensional rational diagrams on $P$: finiteness supplies
a common bound on their degrees.

For completeness, a bounded projective complex has acyclic Hom complex
into every acyclic diagram complex. This follows by induction on its
number of nonzero terms, starting with a shifted projective and using
the short exact sequence separating its top term. Consequently it
computes morphisms in localization by quasi-isomorphisms: the cone of
each inverted map is acyclic, so applying Hom gives a bijection on
chain-homotopy classes; roofs then have unique such representatives.
The desired morphisms are therefore chain-homotopy classes of maps
$Q_P\to B_P$.

The colimit of each representable summand is one copy of $\Q$, and
its boundary maps are exactly deletion of chain vertices. Hence
$\colim_P Q_P=C_*(NP;\Q)$, with its actual augmentation. The
colimit--constant-diagram adjunction, on chain maps and on chain
homotopies, now identifies the preceding classes with
$H^2(\operatorname{Hom}(C_*(NP;\Q),\Q))=H^2(BP;\Q)$.
This is a calculation in the strict diagram localization itself;
it assumes no general comparison with a path-space model.
\end{proof}

\begin{proposition}[C.26: the loop formulas do not imply global invariance]
\label{F:C26-complexes}
The countable category of bounded finite-dimensional rational chain
complexes, marked at quasi-isomorphisms, satisfies both pointed loop and
suspension formulas but fails the global invariance condition of IX.
\end{proposition}
\begin{proof}
\textbf{The precise model hypotheses.} Use finite rational matrices to
choose a countable small model, with degreewise injections as
cofibrations and degreewise surjections as fibrations. It has all finite
limits and colimits, the exact hypothesis M0 of Dwyer--Kan
\cite[Section~1.4(iv)]{dwyerkan1980function}; arbitrary coproducts are
not required by that source. An acyclic vector-space complex is
contractible by splitting its cycles. In a lifting square, a graded lift
has a differential defect in the Hom complex from the cokernel of the
injection to the kernel of the surjection. If either given map is a
quasi-isomorphism, that Hom complex is contractible, and subtracting a
boundary corrects the defect. This proves the lifting axiom.

\textbf{Factorizations and base change.} Write $C(A)=\operatorname{Cone}(1_A)$.
The maps $A\to B\oplus C(A)\to B$, the first $(f,j)$ and the second
projection, are a cofibration followed by a trivial fibration. The maps
$A\to A\oplus C(B)[-1]\to B$, the first inclusion and the second $(f,q)$,
are a trivial cofibration followed by a fibration. Kernels, cokernels,
and exactness of vector spaces verify all pullback and pushout axioms;
homology gives two-out-of-three. Every object is cofibrant and fibrant.

\textbf{The mapping spaces.} The cosimplicial complex
$A\otimes N\Q\Delta^\bullet$ augments by quasi-isomorphisms, and its
latching inclusions are the boundary inclusions, hence degreewise
injections. Dwyer--Kan Proposition~4.4 applies, with fibrant target $B$,
and identifies the derived mapping space with the simplicial abelian
group $r\mapsto\Hom_{\mathrm{Ch}}(A\otimes N\Q\Delta^r,B)$.
Its normalized complex is the good truncation
$\tau_{\geq0}\operatorname{Hom}_*(A,B)$: the value on the top simplex is
the normalized generator and its boundary is the alternating face sum;
in degree zero one takes cycles. Thus its $i$th homotopy group is
$H_i\operatorname{Hom}_*(A,B)$. This verifies the actual comparison used
in C.26, including its finite-limit hypothesis.

\textbf{The formulas and their failure to imply invariance.} Shifts
$\Sigma A=A[1]$ and $\Omega B=B[-1]$ give the based loops of these
mapping spaces. Equivalently the displayed cone factorizations compute
$0\amalg_A^h0$ and $0\times_B^h0$; mapping-space universal properties
supply their coherent group laws. For $A=\Q[0]$ and $B=\Q[2]$ the mapping
space is $K(\Q,2)$. Proposition~\ref{F:C20-paths} represents the actual
constant-endpoint diagram-localization functor by the indicated
ind-path object, on all small tests. If it were globally invariant,
Theorem~\ref{F:global-ind-groupoids} would force its value on a face
poset of $S^2$ to equal its value on a point, since both have terminal
fundamental groupoid. Lemma~\ref{F:diagram-sphere} computes those
actual values as $\Q$ and a singleton. This is impossible. This repairs
the dependency on the unavailable C.3/C.7 path-space proof in the
supplied final paragraph, while confirming the claimed counterexample
and every model-category calculation that precedes it.
\end{proof}
